\section{回避}
\begin{enumerate}
  \item 期待保持中立。影响公正审判。
  \item 侦检审书翻鉴。侦查人员、检察人员、审判人员（法官助理）、书记员、翻译、鉴定人（专门知识的人）。
  \item 回避的理由。利害关系、请客送礼、参前不参后。
        \begin{enumerate}
          \item 利害关系。可能影响公正处理的。近亲属扩大解释，祖宗十八代都属于回避的理由。考特殊情形的回避。
                \al{甲一审合议庭审判员，后升迁为了一审法院的审委会委员。案件被发回重审}。
          \item 请客送礼。凭空污人清白 \yjt 应当提供证据。
          \item 参前不参后。合议庭组成人员 \hz 独任审判员。发回重组、再回不限。另开特别程序不算前后。
                \al{甲为一审书记员，案件被发回重审，甲不能再担任本案书记员 \cuo{书记员不是合议庭成员}}。
          \item 无理的申请。当庭驳回，永世不得复议。
        \end{enumerate}
  \item 种类。回避前提是参与了案件。自行回避、申请回避、指令回避。我国没有整体回避。
        \al{甲宣称 A 市中级法院院长乙为当地黑社会“保护伞”，院长应自行回避 \cuo}。
  \item 程序。注意区分是要求哪方回避。\al{被告人甲要求\xhx{公诉人的书记员}乙回避 \yjt 检察长决定}。
        \begin{enumerate}
          \item 当法诉辩书面 \hz 口头申请。不能停止侦查（有紧迫性）。
          \item 决定。一般找老大，老大找组织；谁聘请谁决定。
          \item 口头 \hz 书面、作出生效、谁不爽谁复议（找做出决定的）。\al{向上级提出复议、上诉 \cuo}。
          \item 回避决定后该人员之前的各种行为效力待定 \yjt 谁决定 \yjt 看回第 (c) 项。
        \end{enumerate}
\end{enumerate}
